% ============================================================
% Notes: Fractional Integration (Fractional Calculus)
% Topics: RL fractional integral, semigroup property, kernels,
%         Laplace transform, basic examples
% Source: Richard Herrmann, Fractional Calculus: An Introduction for Physicists
% ============================================================

\section{Fractional calculus: fractional integration}

\subsection{Motivation}
Fractional integration generalizes the $n$-fold integral
\[
I^n f(t)=\underbrace{\int_0^t\int_0^{\tau_1}\cdots\int_0^{\tau_{n-1}}}_{n\ \text{times}}
f(\tau_n)\,d\tau_n\cdots d\tau_1
\]
to non-integer order $\alpha>0$:
\[
I^n \ \longrightarrow\ I^\alpha.
\]
Fractional integrals are \textbf{nonlocal operators}: $(I^\alpha f)(t)$ depends on all
values of $f(\tau)$ for $0\le \tau\le t$. :contentReference[oaicite:1]{index=1}

% ------------------------------------------------------------

\section{Riemann--Liouville (RL) fractional integral}

\subsection{Definition}
For $\alpha>0$, the (left-sided) RL fractional integral is
\begin{equation}
\boxed{
(I_{0+}^\alpha f)(t)
=
\frac{1}{\Gamma(\alpha)}
\int_0^t (t-\tau)^{\alpha-1}f(\tau)\,d\tau,
\qquad t>0.
}
\end{equation}
The kernel is a power law:
\[
M_\alpha(t-\tau)=\frac{(t-\tau)^{\alpha-1}}{\Gamma(\alpha)}.
\]
Thus $I^\alpha$ is a Volterra-type convolution operator.

\subsection{Integer-order consistency}
If $\alpha=n\in\mathbb{N}$, then
\begin{equation}
(I^n f)(t)
=
\frac{1}{(n-1)!}\int_0^t (t-\tau)^{n-1}f(\tau)\,d\tau,
\end{equation}
which reproduces the standard $n$-fold integral.

\subsection{Semigroup property}
For suitable $f$,
\begin{equation}
\boxed{
I^\alpha I^\beta f = I^{\alpha+\beta}f,
\qquad \alpha,\beta>0.
}
\end{equation}

\subsection{Identity element}
\begin{equation}
I^0 f = f.
\end{equation}

% ------------------------------------------------------------

\section{Right-sided RL fractional integral}

On an interval $[a,b]$, the right-sided RL fractional integral is
\begin{equation}
\boxed{
(I_{b-}^\alpha f)(t)
=
\frac{1}{\Gamma(\alpha)}
\int_t^b (\tau-t)^{\alpha-1}f(\tau)\,d\tau.
}
\end{equation}
This is useful in boundary-value problems and symmetric formulations.

% ------------------------------------------------------------

\section{Fractional integral as convolution}

\subsection{Convolution form}
Define
\[
g_\alpha(t)=\frac{t^{\alpha-1}}{\Gamma(\alpha)},\qquad t>0.
\]
Then
\begin{equation}
(I^\alpha f)(t) = (g_\alpha * f)(t)
:= \int_0^t g_\alpha(t-\tau)f(\tau)\,d\tau.
\end{equation}

\subsection{Laplace transform rule}
Let $F(s)=\mathcal{L}\{f(t)\}(s)$. Then
\begin{equation}
\boxed{
\mathcal{L}\{I^\alpha f\}(s) = s^{-\alpha}F(s).
}
\end{equation}
This is one of the most useful operational formulas.

% ------------------------------------------------------------

\section{Examples}

\subsection{Fractional integral of a constant}
Let $f(t)=1$. Then
\begin{equation}
(I^\alpha 1)(t)
=
\frac{1}{\Gamma(\alpha)}
\int_0^t (t-\tau)^{\alpha-1}\,d\tau
=
\frac{t^\alpha}{\Gamma(\alpha+1)}.
\end{equation}

\subsection{Fractional integral of a power}
Let $f(t)=t^\mu$ with $\mu>-1$. Then
\begin{equation}
\boxed{
(I^\alpha t^\mu)(t)
=
\frac{\Gamma(\mu+1)}{\Gamma(\mu+\alpha+1)}\,t^{\mu+\alpha}.
}
\end{equation}

\subsection{Fractional integral of an exponential (series form)}
Using $e^{at}=\sum_{n=0}^\infty \frac{(at)^n}{n!}$ and linearity:
\begin{equation}
I^\alpha e^{at}
=
\sum_{n=0}^\infty \frac{a^n}{n!}\,I^\alpha t^n
=
\sum_{n=0}^\infty \frac{a^n}{n!}\frac{\Gamma(n+1)}{\Gamma(n+\alpha+1)}t^{n+\alpha}.
\end{equation}

% ------------------------------------------------------------

\section{Important remarks}

\subsection{Nonlocality / memory}
Because
\[
(I^\alpha f)(t)=\frac{1}{\Gamma(\alpha)}\int_0^t (t-\tau)^{\alpha-1}f(\tau)\,d\tau,
\]
the value at time $t$ depends on the full past history $0\le\tau\le t$ with a power-law weight.
This is the origin of memory effects in fractional models. :contentReference[oaicite:2]{index=2}

\subsection{Kernel behavior}
For $0<\alpha<1$, the kernel behaves like $(t-\tau)^{\alpha-1}$ which is integrable but singular
near $\tau=t$:
\[
(t-\tau)^{\alpha-1}\to \infty \quad \text{as }\tau\to t^-.
\]
So recent history has the strongest weight, but all past contributes.

% ------------------------------------------------------------

\section{High-yield summary}

\begin{itemize}
\item RL fractional integral ($\alpha>0$):
\[
(I^\alpha f)(t)=\frac{1}{\Gamma(\alpha)}\int_0^t (t-\tau)^{\alpha-1}f(\tau)\,d\tau.
\]
\item Semigroup property:
\[
I^\alpha I^\beta = I^{\alpha+\beta}.
\]
\item Laplace transform:
\[
\mathcal{L}\{I^\alpha f\}=s^{-\alpha}F(s).
\]
\item Examples:
\[
I^\alpha 1 = \frac{t^\alpha}{\Gamma(\alpha+1)},\qquad
I^\alpha t^\mu = \frac{\Gamma(\mu+1)}{\Gamma(\mu+\alpha+1)}t^{\mu+\alpha}.
\]
\end{itemize}
